\section{Introduction}
\label{sec:intro}

Machine-learned models increasingly mediate consequential decisions about
people: whether a defendant is labelled high-risk, whether an applicant
receives credit, whether a r\'esum\'e is advanced. A decade of work has
established that such models can encode and amplify disparities present in
their training data~\cite{angwin2016machine,buolamwini2018gender,mehrabi2021survey}.
What is far less settled is what should be done about it---because
``fairness'' is not one property but a family of mutually inconsistent ones.

Two families dominate practice. \emph{Group fairness} requires that outcomes
be comparable across demographic groups; its canonical instance, demographic
parity, requires equal selection rates. \emph{Individual fairness} requires
that similar individuals be treated similarly~\cite{dwork2012fairness}, and in
its influential outcome-based operationalization requires that the
\emph{benefit} each person derives from the classifier be distributed evenly
across the population~\cite{speicher2018unified}. These are different
requirements, and satisfying one does not imply satisfying the other.

That they can conflict is known theoretically. Calibration and error-rate
balance cannot generally hold at once when base rates
differ~\cite{kleinberg2016inherent,chouldechova2017fair}; demographic parity
and equalized odds are likewise incompatible under the same
condition~\cite{barocas2019fairness}. What these results do \emph{not} supply
is a magnitude. A practitioner told that ``you cannot have both'' still needs
to know what one costs in units of the other, on data resembling their own,
with uncertainty quantified. Comparative empirical studies of fairness
interventions exist~\cite{friedler2019comparative}, but they typically compare
\emph{discrete methods}, each committed in advance to a single fairness
notion. What is missing is a study that varies the fairness objective
\emph{continuously} and measures both families of metrics along the way.

\subsection{Contributions}

We build such an instrument and use it to map the trade-off.

\begin{enumerate}
\item \textbf{A differentiable composite loss} (Section~\ref{sec:method}) whose
      parameter $\beta$ interpolates continuously between a soft
      demographic-parity term and a soft generalized-entropy term, with a
      second parameter $\alpha$ controlling the overall weight on fairness
      versus accuracy. Both fairness terms are computed on raw predicted
      probabilities so that gradients flow through them---the thresholding
      step that defines the evaluation metrics is not differentiable and
      cannot be optimized directly.
\item \textbf{An empirical trade-off map} (Section~\ref{sec:results}) built
      from a $7\times7$ $(\alpha,\beta)$ grid across three standard datasets,
      two model families, and ten seeds: $2{,}940$ trained models, all metrics
      reported as mean~$\pm$~standard deviation with paired significance tests.
\item \textbf{A statistically significant double dissociation.} Optimizing the
      group term reduces the demographic parity difference by $78\%$--$89\%$
      ($p_{\text{Holm}}=0.008$) while significantly \emph{worsening}
      individual-fairness inequality; optimizing the individual term
      significantly improves individual fairness while leaving group disparity
      unchanged or significantly worse. On Adult, driving demographic parity
      to near zero degrades equalized odds by a factor of $3.3$, an empirical
      instance of a known impossibility.
\item \textbf{A diagnosed negative result.} On German Credit no effect
      survives correction. We show this is not a failure of the method but of
      \emph{validation-based model selection at small $n$}: validation and
      test demographic parity are effectively uncorrelated
      ($r=0.08$ versus $0.94$ and $0.99$ elsewhere), so the selection rule
      cannot identify fair configurations. Under a pre-specified
      configuration the effect is present. We believe this sample-size floor
      on fairness model selection is of independent practical interest.
\item \textbf{A benchmark soundness correction.} We identify and remove a
      label-leaking feature from a COMPAS feature set in prior use, and report
      the honest accuracy that results.
\end{enumerate}

The composite loss is the \emph{instrument}; the trade-off map is the
\emph{contribution}. We do not claim a new state of the art in fairness
mitigation, and Section~\ref{sec:results} shows the loss matches rather than
beats established baselines on group-fairness metrics. Its value is that,
unlike those baselines, it exposes a continuous dial between two fairness
notions, which is precisely what mapping a trade-off requires.
